% !TEX program = xelatex
% !TEX root = ../../TFM.tex
\documentclass[../../TFM.tex]{subfiles}

\begin{document}

	% =====================================================
	% CRONOGRAMA
	% =====================================================

	\chapter{Cronograma}
	\label{chap:cronograma}

	La \autoref{tab:cronograma} resume el cronograma tentativo de las actividades previstas para el desarrollo del trabajo durante el año 2026, en coherencia con las etapas metodológicas descritas en el capítulo metodológico del proyecto.

	\begin{longtable}{|p{6.2cm}|*{12}{c|}}
		\caption[Cronograma de actividades]{Cronograma tentativo de las actividades de investigación durante el año 2026.}

		\label{tab:cronograma} \\
		\hline
		\textbf{Actividad} & \multicolumn{12}{c|}{\textbf{Meses del año 2026}} \\
		\hline
		& 1 & 2 & 3 & 4 & 5 & 6 & 7 & 8 & 9 & 10 & 11 & 12 \\
		\hline
		\endfirsthead

		\multicolumn{13}{c}%
		{{\tablename\ \thetable{} -- continuación}} \\
		\hline
		\textbf{Actividad} & \multicolumn{12}{c|}{\textbf{Meses del año 2026}} \\
		\hline
		& 1 & 2 & 3 & 4 & 5 & 6 & 7 & 8 & 9 & 10 & 11 & 12 \\
		\hline
		\endhead

		\hline
		\multicolumn{13}{r}{\textit{Continúa en la página siguiente}} \\
		\endfoot

		\hline
		\endlastfoot

		\textbf{1.} Revisión exploratoria de literatura y organización bibliográfica (artículos académicos, Mendeley)
		& \cellcolor{gray!40} & \cellcolor{gray!40} &  &  &  &  &  &  &  &  &  &  \\
		\cline{2-13}

		\textbf{2.} Ajustes del marco teórico y revisión dirigida de literatura
		&  & \cellcolor{gray!40} & \cellcolor{gray!40} & \cellcolor{gray!40} &  &  &  &  &  &  &  &  \\
		\cline{2-13}

		\textbf{3.} Construcción, extracción y limpieza de la base de datos
		&  & \cellcolor{gray!40} & \cellcolor{gray!40} & \cellcolor{gray!40} &  &  &  &  &  &  &  &  \\
		\cline{2-13}

		\textbf{4.} Especificación formal del modelo econométrico y definición final de variables
		&  &  & \cellcolor{gray!40} & \cellcolor{gray!40} &  &  &  &  &  &  &  &  \\
		\cline{2-13}

		\textbf{5.} Análisis descriptivo y caracterización de trayectorias
		&  &  & \cellcolor{gray!40} & \cellcolor{gray!40} & \cellcolor{gray!40} &  &  &  &  &  &  &  \\
		\cline{2-13}

		\textbf{6.} Estimación principal con ECI y complementaria con DIVX
		&  &  &  & \cellcolor{gray!40} & \cellcolor{gray!40} & \cellcolor{gray!40} &  &  &  &  &  &  \\
		\cline{2-13}

		\textbf{7.} Análisis de resultados e integración con el marco analítico
		&  &  &  &  &  & \cellcolor{gray!40} & \cellcolor{gray!40} &  &  &  &  &  \\
		\cline{2-13}

		\textbf{8.} Redacción de resultados
		&  &  &  &  &  & \cellcolor{gray!40} & \cellcolor{gray!40} & \cellcolor{gray!40} &  &  &  &  \\
		\cline{2-13}

		\textbf{9.} Redacción de conclusiones y recomendaciones
		&  &  &  &  &  &  &  & \cellcolor{gray!40} & \cellcolor{gray!40} &  &  &  \\
		\cline{2-13}

		\textbf{10.} Correcciones finales, revisión integral y entrega
		&  &  &  &  &  &  &  &  & \cellcolor{gray!40} & \cellcolor{gray!40} & \cellcolor{gray!40} &  \\

		\hline
		\multicolumn{13}{p{\textwidth}}{%
			\textit{Nota: El cronograma es tentativo y puede ajustarse según el avance del trabajo y la disponibilidad de información.}
		}

	\end{longtable}

\end{document}
